
\documentclass[11pt,a4paper]{article}
\usepackage[numbers]{natbib}
\usepackage{booktabs}
\usepackage{hyperref}
\usepackage{tikz}
\usepackage{amsmath}
\usepackage{geometry}
\geometry{margin=1in}

\title{machine learning: A Survey of Recent arXiv Research}
\author{Redigg AI Research}
\date{\today}

\begin{document}
\maketitle

\begin{abstract}
This survey reviews recent advances in machine learning based on analysis of 10 papers from arXiv (2026-04-02 to 2026-04-02). We identify key trends including gpt, llm, vit. Our analysis reveals significant progress with applications spanning multiple domains. We discuss open challenges and outline promising directions for future research.
\end{abstract}


\section{Introduction}

Recent advances in machine learning have attracted significant attention. This survey analyzes 10 recent papers from arXiv.

\subsection{Motivation}

The rapid growth of publications makes it challenging to track developments. This survey synthesizes key findings.

\subsection{Contributions}

\begin{itemize}
    \item Analysis of 10 recent arXiv papers
    \item Identification of 10 key themes
    \item Taxonomy of methods
    \item Open challenges and future directions
\end{itemize}



\section{Background}

\subsection{Historical Context}

Research in machine learning has evolved significantly.

\subsection{Key Concepts}

Foundational concepts include attention mechanisms, representation learning, and transfer learning.



\section{Taxonomy}

We categorize recent work:

\begin{itemize}
    \item \textbf{Gpt}: Active research area
    \item \textbf{Llm}: Active research area
    \item \textbf{Vit}: Active research area
    \item \textbf{Diffusion}: Active research area
    \item \textbf{Graph}: Active research area
\end{itemize}



\section{Methodologies}

\begin{itemize}
    \item \textbf{Gpt}: Emerging technique
    \item \textbf{Llm}: Emerging technique
    \item \textbf{Vit}: Emerging technique
    \item \textbf{Diffusion}: Emerging technique
\end{itemize}



\section{{Experimental Analysis}}

\subsection{{Benchmarks}}

Papers evaluate on standard benchmarks.

\subsection{{Metrics}}

Common metrics include accuracy, efficiency, and generalization.



\section{{Challenges}}

\subsection{{Scalability}}

Large-scale scenarios remain challenging.

\subsection{{Efficiency}}

Computational efficiency is an active area.



\section{{Future Directions}}

\subsection{{Emerging Trends}}

Continued advancement in architectures and methods.

\subsection{{Applications}}

Expanding applications to new domains.



\section{Conclusion}

This survey reviewed recent advances in machine learning. The field continues to evolve rapidly.


\begin{thebibliography}{99}

\bibitem{Ref1} Alexander Pondaven and Ziyi Wu and Igor Gilitschenski and others. "ActionParty: Multi-Subject Action Binding in Generative Video Games." arXiv:2604.02330v1, 2026.
\bibitem{Ref2} Alex Costanzino and Pierluigi Zama Ramirez and Giuseppe Lisanti and others. "Modulate-and-Map: Crossmodal Feature Mapping with Cross-View Modulation for 3D Anomaly Detection." arXiv:2604.02328v1, 2026.
\bibitem{Ref3} Daiwei Chen and Zhoutong Fu and Chengming Jiang and others. "Grounded Token Initialization for New Vocabulary in LMs for Generative Recommendation." arXiv:2604.02324v1, 2026.
\bibitem{Ref4} Ruozhen He and Nisarg A. Shah and Qihua Dong and others. "Beyond Referring Expressions: Scenario Comprehension Visual Grounding." arXiv:2604.02323v1, 2026.
\bibitem{Ref5} Bangji Yang and Hongbo Ma and Jiajun Fan and others. "Batched Contextual Reinforcement: A Task-Scaling Law for Efficient Reasoning." arXiv:2604.02322v1, 2026.
\bibitem{Ref6} Junxuan Li and Rawal Khirodkar and Chengan He and others. "Large-scale Codec Avatars: The Unreasonable Effectiveness of Large-scale Avatar Pretraining." arXiv:2604.02320v1, 2026.
\bibitem{Ref7} Yuhan Liu and Fangyuan Xu and Vishakh Padmakumar and others. "No Single Best Model for Diversity: Learning a Router for Sample Diversity." arXiv:2604.02319v1, 2026.
\bibitem{Ref8} Christian Ferko and James Halverson and Vishnu Jejjala and others. "Topological Effects in Neural Network Field Theory." arXiv:2604.02313v1, 2026.
\bibitem{Ref9} Torque Dandachi and Sophia Diggs-Galligan. "go-$m$HC: Direct Parameterization of Manifold-Constrained Hyper-Connections via Generalized Orthostochastic Matrices." arXiv:2604.02309v1, 2026.
\bibitem{Ref10} Prakul Sunil Hiremath and PeerAhammad M Bagawan and Sahil Bhekane. "PARD-SSM: Probabilistic Cyber-Attack Regime Detection via Variational Switching State-Space Models." arXiv:2604.02299v1, 2026.

\end{thebibliography}

\end{document}
